\documentclass[11pt]{article}
\usepackage{amsmath,amssymb,amsthm,booktabs,hyperref}
\newtheorem{theorem}{Theorem}
\newtheorem{lemma}[theorem]{Lemma}
\newtheorem{proposition}[theorem]{Proposition}
\title{Finite enumeration of the right halves in the singleton middle case}
\author{human\_compression project}
\date{}
\begin{document}
\maketitle

\paragraph{Scope.}
The retained singleton-middle theorem~\cite{middle} reduces its ordinary
graph stratum to six halves of order fourteen and a class of halves of
order sixteen. The latter class has cardinality $h_{16}$, still
undetermined. This note specifies an exhaustive finite parameter domain
and an exact ordinary canonical code for that class. No enumeration,
control, graph construction, TeX compilation or Lean execution was
performed by its author. In particular no numerical value of $h_{16}$
or of the original graph census is asserted.

The degree and neighborhood reductions, positive incidence condition,
automatic diameter theorem and product $6h_{16}$ are borrowed from the
retained theorem. The original cage problem belongs to Cambie,
Goedgebeur, Jooken and Van den Eede~\cite{cages}. Independent-row margins,
multisets, automorphisms of small forests and intrinsic vertex degrees
are standard ingredients. The new project assembly is the finite domain,
its ordinary identification and its bounded execution interface; no
historical priority is claimed.

\section*{The finite parameter domain}

A right half has a distinguished endpoint $p$, five endpoint rows $A$,
seven side vertices $D$, and three ports $P$. Its internal side graph is
one of
\[
 D_1=P_4\sqcup3K_1,\quad D_2=K_{1,3}\sqcup3K_1,\quad
 D_3=P_3\sqcup K_2\sqcup2K_1,\quad D_4=3K_2\sqcup K_1.
\]
Every endpoint row meets $p$ and an independent four-subset of $D$.
Every port meets an independent three-subset of $D$. There are no other
edges except the three edges of $D$. Repeated row subsets denote
distinct row vertices and are permitted.

Write $\mathcal S$ for the multiset of five four-subsets and
$\mathcal T$ for the multiset of three three-subsets. For a side vertex
$v$, let $\alpha(v)$ and $\beta(v)$ be its multiplicities in these two
multisets. The complete feasibility condition is
\begin{equation}\label{margins}
 \alpha(v)+\beta(v)+d_D(v)=5,\qquad \alpha(v)\geq1.
\end{equation}
Positivity is essential: a zero endpoint column would put that vertex
in the middle and invalidate the singleton presentation.

\begin{proposition}\label{bounds}
The numbers of permitted individual row types and of their unrestricted
multisets are as follows:
\[
\begin{array}{c|rrrr|r}
 D&a=\#\mathcal I_4(D)&b=\#\mathcal I_3(D)
  &N_A=\binom{a+4}{5}&N_B=\binom{b+2}{3}&|\operatorname{Aut}D|\\\hline
 D_1&13&22&6188&2024&12\\
 D_2&16&23&15504&2300&36\\
 D_3&11&21&3003&1771&8\\
 D_4&8&20&792&1540&48
\end{array}
\]
The sum of the four raw Cartesian bounds $N_A N_B$ is $54\,721\,705$.
These are parameter bounds, not graph counts or workload benchmarks.
\end{proposition}
\begin{proof}
The independence polynomials of the four side graphs are
\[
\begin{split}
 &(1+4z+3z^2)(1+z)^3,\\
 &((1+z)^3+z)(1+z)^3,\\
 &(1+3z+z^2)(1+2z)(1+z)^2,\\
 &(1+2z)^3(1+z).
\end{split}
\]
Their coefficients of $z^4,z^3$ give $a,b$. Multisets give the binomial
counts displayed. For $D_1$, path reversal and permutation of the three
isolates give $2\cdot6=12$ automorphisms. For $D_2$, permutations of the
three leaves and the three isolates give $6\cdot6=36$. For $D_3$ the
two path ends, two edge ends and two isolates give $2^3=8$; their
components are distinguishable. For $D_4$ there are $3!$ edge orders and
$2^3$ choices of orientations, giving $48$. Component structure shows
that these lists are exhaustive.
\end{proof}

It is unnecessary to visit the entire raw Cartesian product. Enumerate
the $N_A$ endpoint multisets once and index them by their seven column
counts. For a port multiset, compute
$\beta$ and the forced target $\alpha(v)=5-d_D(v)-\beta(v)$.
If some target entry is nonpositive there is no completion; otherwise
consult only the matching endpoint list. The largest index has at most
$15\,504$ records, each carrying a multiset of five row types. This is a
finite structural reduction, not a prediction of a measured memory fit.

\section*{Removing labels of the marked presentation}

Let $\Gamma=\operatorname{Aut}(D)$. Row labels have already been removed
by using multisets. First require $\mathcal T$ to be least under
$\Gamma$. Its stabilizer $\Gamma_{\mathcal T}$ acts on the matching
endpoint multisets. Retain $\mathcal S$ only when it is least under
that stabilizer.

\begin{lemma}\label{marked}
This procedure retains exactly one representative of each marked
incidence object, allowing permutations of the side graph, endpoint rows
and ports, while keeping the endpoint $p$ distinguished.
\end{lemma}
\begin{proof}
An isomorphism fixing the marked endpoint preserves its endpoint rows,
the degree-three ports, and the remaining side vertices. Its side
restriction belongs to $\Gamma$ and acts simultaneously on both row
multisets. Every orbit has a least port multiset. Once that multiset is
fixed, all remaining choices differ precisely by its stabilizer. That
stabilizer preserves $\beta$ and $d_D$, and hence the target $\alpha$;
it therefore acts on the whole matching endpoint domain. Taking its
minimum gives exactly one representative. Conversely each such action
lifts to the indicated row and side bijections.
\end{proof}

Neither these marked representatives nor separate counts for the four
side graphs are ordinary half classes. The distinguished endpoint may
change under an unrooted isomorphism, and the induced side graph may
then have a different one of the four types.

\section*{Ordinary identification without a preferred root}

In an admissible half $H$, exactly three vertices have degree three;
they form the intrinsic port set $P$. The other thirteen have degree
five. The half is connected, simple and has no triangles. All these
properties are directly checked in the proposed implementation.

An eligible endpoint is a degree-five vertex $p$ for which, on putting
$A=N_H(p)$ and $D=V(H)\setminus(\{p\}\cup A\cup P)$,
\begin{enumerate}
\item $A$ contains no port and has no edge to a port;
\item every vertex of $D$ has a neighbor in $A$.
\end{enumerate}
The ports themselves are independent. There are at most thirteen
eligible endpoints, and at least the original one exists.

\begin{lemma}\label{roots}
These conditions characterize exactly the endpoint choices admitting
the required right-half presentation. For every such choice, $D$ has
seven vertices and three edges and is one of $D_1,\ldots,D_4$.
\end{lemma}
\begin{proof}
They are necessary in every presentation. Conversely $A$ is independent
because it is the neighborhood of $p$ in a graph without triangles.
Its vertices have $p$ and four neighbors in $D$: the other possibilities
are excluded by the two port conditions and the definition of $A$.
Every port has all three neighbors in $D$. Thus the rows have exactly
the required sizes and are independent. Their column incidences meet
\eqref{margins}, with positivity by the second condition.

The half has $(13\cdot5+3\cdot3)/2=37$ edges. Five join $p$ to $A$,
twenty join $A$ to $D$, and nine join ports to $D$. The remaining three
lie in $D$. The graph $D$ has no triangles. A graph with three edges
and no triangles is either $P_4$, $K_{1,3}$, $P_3\sqcup K_2$, or
$3K_2$, padded with isolates to order seven. This proves the claim and
recovers every adjacency of the presentation.
\end{proof}

Fix a standard labeling of each $D_i$. For each eligible $p$, inspect
every isomorphism from its side graph to that labeling. There are
$|\operatorname{Aut}D_i|\leq48$ such maps. They can be listed from
components: choose a path orientation, leaf order, edge orientation or
edge-component order as appropriate, and permute isolates. For each map
record the endpoint and port neighbor subsets as seven-bit masks, sort
the five endpoint masks and three port masks, and form
\[
 (i,\operatorname{sort}\mathcal T,\operatorname{sort}\mathcal S).
\]
Let $C(H)$ be the least of these tuples over all eligible endpoints.
There are at most $13\cdot48=624$ candidate tuples. Sorting removes
row labels; no additional port permutation or preferred endpoint is kept.

\begin{theorem}\label{ordinary}
For admissible halves $H,H'$, equality $C(H)=C(H')$ holds if and only
if they are isomorphic as ordinary unrooted graphs on sixteen vertices.
\end{theorem}
\begin{proof}
An ordinary isomorphism preserves vertex degrees and hence $P$. It maps
eligible endpoints to eligible endpoints by their intrinsic conditions,
and maps the associated $A,D$ sets and every side isomorphism to the
corresponding choices. The two complete sets of tuples are therefore
equal, so their minima agree.

Conversely a tuple reconstructs the entire half. Allocate its endpoint,
five endpoint rows, seven standard side vertices and three ports in
that order; add the three side edges, the five endpoint edges, and the
two sorted row multisets. Equality of tuples gives a common complete
graph. The maps attaining the tuples, with any order among repeated
equal row subsets, give explicit vertex bijections to that graph.
Composing them gives an ordinary isomorphism. Thus the tuple is a
complete invariant, rather than a partial invariant or an unchecked hash.
\end{proof}

The proposed program retains the sixteen-vertex bijection attaining
each code and verifies the entire adjacency relation against its
reconstruction. This certifies every claimed identification directly.
The completeness of the canonicalizer follows from the endpoint and
side-map arguments above, not merely from agreement on test fixtures.

\section*{Bounded intervals and the final count}

For each side type give its endpoint multisets raw ordinals
$i\in[0,N_A)$ and its port multisets ordinals $j\in[0,N_B)$ in fixed
lexicographic orders. The raw cell ordinal is $jN_A+i$. A job covers a
specified half-open interval in this finite rectangle. Its endpoint
index proves which cells have incompatible column margins; the
automorphism argument proves which blocks or cells are aliases. Such
cells may be skipped in a bounded step without constructing graphs.

Every compatible representative is constructed and receives the full
ordinary code and permutation witness. A job advances its frontier
only past cells whose rejection, alias or complete output has been
established. A stop in a compatible cell leaves that entire cell next
and unpaid. Data are streamed with explicit record and byte limits;
the clock, RAM, disk, declared-input and ownership guards remain
separate requirements. No expected literature count is a stopping rule.

\begin{theorem}\label{coverage}
After actual reviewed jobs cover every raw cell for all four side types,
the distinct ordinary codes in their combined output are in bijection
with $\mathcal H_{16}$.
\end{theorem}
\begin{proof}
The finite domain contains every possible endpoint and port row
multiset. Indexing by margins removes only violations of
\eqref{margins}; the two automorphism minima remove only marked
isomorphic copies, retaining a representative of every such orbit.
Every retained object constructs a feasible half, and every feasible
half has a marked presentation somewhere in this exhaustive domain.
Theorem~\ref{ordinary} identifies precisely its ordinary copies,
including copies arising from different endpoint choices and side types.
Taking the complete union therefore gives exactly the stated bijection.
\end{proof}

This theorem does not assign a numerical value to $h_{16}$. Coverage,
actual output, independent review and global ordinary identification
remain owed. The formula $6h_{16}$ supplies only the corresponding
singleton stratum of the original cage problem. It cannot determine
$h_{16}$ from an unrelated total for all nonbipartite graphs, and it
does not pay an existing native queue record without a separate exact
family/coverage application.

\begin{thebibliography}{9}
\bibitem{middle} Human Compression project,
\emph{The empty and singleton middle cases for quintic graphs}, frozen
proof source SHA256
\texttt{76a5c09c2e3f52536c022572bd95ba611392ad38527c0f2ee40c1e886e9d9eb2},
with its separate independent mathematical review.
\bibitem{cages} Cambie, Goedgebeur, Jooken and Van den Eede,
\emph{On the order-diameter ratio of girth-diameter cages},
arXiv:2511.21144v3, Appendix~B, Table~6.
\end{thebibliography}
\end{document}
